\uuid{kwhO}
\exo7id{829}
\titre{exo7 829}
\auteur{ridde}
\organisation{exo7}
\datecreate{1999-11-01}
\isIndication{false}
\isCorrection{false}
\chapitre{Calcul d'intégrales}
\sousChapitre{Fraction rationnelle en sin, cos ou en sh, ch}
\cpge{mpsi-mp2i.primitives-equations-differentielles.a, pcsi.techniques-calcul-primitives.a, ptsi.techniques-calcul-primitives.a, mpsi-mp2i.fonctions-valeurs-reelles-complexes.b, pcsi.techniques-calcul-fonctions.b, ptsi.techniques-calcul-fonctions.b, mpsi-mp2i.calcul-algebrique-trigonometrie.d, pcsi.calcul-algebrique-trigonometrie.d, ptsi.calcul-algebrique-trigonometrie.d, bcpst1.derivees-primitives-integrales.c, bcpst1.derivees-primitives-integrales.d, bcpst1.trigonometrie}
\module{Analyse}
\niveau{L1}
\difficulte{}

\contenu{
\question{
D\'eterminer les intervalles d'\'etude et calculer les primitives des fonctions :
$$\frac{\cos^3 x}{\sin x}$$
$$ \frac 1{1 + \tan x}$$
$$\frac 1 {\text{th}^2 x}$$
}
}
